\section{Physics of the 21cm Signal}
For the purposes of this thesis, a flat $\Lambda$CDM cosmology with an additional warm dark matter parameter is assumed. The cosmological setting is that of the matter-dominated universe after recombination, during which the CMB was formed. During this time, most photons propagate freely and are quickly redshifted, such that the universe becomes dark. This period is therefore called the Dark Ages. Matter has not yet coalesced into stars and galaxies. Once first stars and galaxies have formed, they emit photons energetic enough to ionize the surrounding hydrogen. This period is called the Epoch of Reionization. By the end of the EoR, most of the IGM is ionized and continues to be so today.

This section covers $21~\si{cm}$ line of neutral hydrogen and how the $21~{cm}$ from this cosmological period has evolved until today. It is based on~\cite{fundamentals}.

\subsection{The origin of the 21cm signal}
The two electron spin states of neutral hydrogen in the $1s$ state have slightly different energy levels. This hyperfine structure originates from the parallel (triplet state) and antiparallel (singlet state) orientation of electron spin with respect to the nucleus. The transition between states is forbidden, the excited state having a mean lifetime of $\sim 11\cdot 10^6~\si{yrs}$. The triplet state carries slightly more energy, the energy difference being $\Delta E \approx 10^{-6}~\si{eV}$, which corresponds to a frequency of $\nu \approx 1.42~\si{GHz}$, or a wavelength of $\lambda \approx 21~\si{cm}$. 

Neutral hydrogen gas made up most of baryonic matter during CD and EoR, and $21~\si{cm}$ photons emitted during these epochs can be measured today. A $21~\si{cm}$ photon emitted at some point during CD or EoR is redshifted by the time it reaches the earth. At redshift ranges $z \sim 6-30$ , it is now measured on a range of $\nu\sim 46-200~\si{MHz}$, or $\lambda\sim 1.5-6.5~\si{m}$. This wavelength range belongs to the regime of radio astronomy and is covered by the SKA-Low setup, which is sensitive to frequencies $\nu \sim (50-350)~\si{MHz}$\footnote{See the \href{https://www.skao.int/en/science-users/118/ska-telescope-specifications}{SKAO website}.}.

\subsection{The evolution of the 21cm signal until today}
There are various mechanisms that lead to changes in this spin state of neutral hydrogen in the IGM. Those mechanisms that are consequences of the quantum mechanical properties of the atom are measured by Einstein coefficients, which are introduced in the following. The singlet state is denoted with index $0$ and the triplet state with index $1$. Then the index $01$ denotes excitation, and $10$ de-excitation. 

$A_{10}$ describes spontaneous emission and can be given as a rate. For neutral hydrogen, one finds $A_{10}\approx 2.85\cdot 10^{-15}~\si{s}^{-1}$.

$B_{01/10}$ describes absorption/stimulated emission. For neutral hydrogen, one finds $B_{10} = A_{10}(c^2 /2h\nu^3)$ and $B_{10}g_1 = B_{01}g_0$, where $\nu$ is the $21~\si{cm}$ line frequency given above and $g_i$ is the spin degeneracy factor of each state. In terms of the $21~\si{cm}$ it is used to describe the interaction of neutral hydrogen clouds with CMB radiation.

Two other effects also lead to the (de-)excitation of neutral hydrogen in the IGM, and their rates are described by effective Einstein coefficients. C denotes those processes driven by collision interactions with other hydrogen atoms, free electrons and protons, and P those from scattering with Lyman-$\alpha$ photons. In equilibrium, both define a kinetic temperature $T_K$ and effective color temperature $T_c$:
\begin{align}
    \frac{C_{01}}{C_{10}} &= \frac{g_1}{g_0} e^{-\frac{T_*}{T_K}}\\
    \frac{P_{01}}{P_{10}} &= \frac{g_1}{g_0} e^{-\frac{T_*}{T_c}},
\end{align}
where again $T_K,T_c \gg T_*$ and the exponential can be expressed in first-order Taylor expansion.

On a large scale, the relative populations $n_0$ and $n_1$ of neutral hydrogen atoms in different spin states are described by (or determine) the spin temperature $T_S$:
\begin{align}
    \frac{n_1}{n_0} 
    = \frac{g_1}{g_0} \exp \biggl(-\frac{E_{10}}{k_B T_S}\biggr) 
    \equiv \frac{g_1}{g_0} \exp \biggl(-\frac{T_\star}{T_S}\biggr),~\label{eq:Boltzmann}
\end{align} 
As a reminder, the $g_i$ are spin degeneracy factors and $E_{10}$ the energy difference between spin states. An estimate of $T_S$ allows an estimate of the relative populations. For realistic astrophysical environments, $T_S \gg T_*$ and as such $n_1 \approx 3n_0$ is a working approximation. 

An estimate of $T_S$ is then obtained by a rate balancing of the aforementioned Einstein coefficients:
\begin{align}
    n_1 (A_{10} + B_{10}I_\gamma +C_{10}+P_{10})=n_0(B_{01}I_\gamma + C_{01}+P_{01}),~\label{eq:rate_balancing}
\end{align}
where $I_\gamma$ is the specific intensity of CMB photons at the frequency of the $21~\si{cm}$ line. The CMB background is the source closest resembling a blackbody spectrum, which is in this frequency range described by the Rayleigh-Jeans law, with which $I_\gamma$ can be written in terms of a temperature at the frequency of the $21~\si{cm}$ line. Plugging everything into~\ref{eq:rate_balancing} yields the following approximation to $T_S$:
\begin{align}
    T_S^{-1} &= \frac{T_\gamma^{-1} + x_c T_K^{-1} + x_\alpha T_C^{-1}}{1+x_c+x_\alpha}\\
    x_c &= \frac{C_{10}}{A_{10}}\frac{h\nu}{k_B T_\gamma}\\
    x_\alpha &= \frac{P_{10}}{A_{10}}\frac{h\nu}{k_B T_\gamma}.
\end{align}
This raises the question of what the measured $21~\si{cm}$ looks like, given this estimate of $T_S$. 

The $21~\si{cm}$ signal is given in terms of brightness offset temperature $\delta T_B$ between the CMB and the spin temperatures and is calculated from the specific intensity measured by the radio telescope, assuming the Rayleigh-Jeans regime. This is done with the radiative transfer equation
\begin{align}
    \frac{dI_\nu}{ds} = -\alpha(\nu) I_\gamma + \alpha(\nu) I_S,~\label{eq:specific_intensity}
\end{align}
where $\nu$ is the frequency, $I_\gamma$ denotes the CMB specific intensity, $I_S$ that of the $21~\si{cm}$ signal from the neutral hydrogen cloud and $\alpha(\nu)$ the absorption coefficient of neutral hydrogen.

The observed specific intensity $I_B$ is obtained by solving the radiative transfer equation from~\ref{eq:specific_intensity}, yielding:
\begin{align}
    I_B &= I_S (1-e^{-\tau})+I_\gamma e^{-\tau}\\
    \tau(\nu) &= \int \alpha(\nu) ds,
\end{align}
where $\tau$ is the so-called optical depth. This carries the interpretation that for an optically thick material, $\tau \gg 1$, the signal is dominated by the $21~\si{cm}$ signal and for an optically thin material,  $\tau \ll 1$, by CMB radiation. 

Once again using the Rayleigh-Jeans law to convert this to temperature and adding the redshift correction for the CMB temperature, this yields
\begin{align}
    T_{B0}(1+z) = T_S (1-e^{-\tau})+T_0 (1+z)e^{-\tau}. 
\end{align}
Finally, the brightness offset temperature is obtained by subtracting the CMB background:
\begin{align}
    \delta T_B = T_{B0} - T_0 = \frac{T_S - T_\gamma}{1+z}(1-e^{-\tau}).~\label{eq:delta_T_simple}
\end{align}
The quantity $\tau$ remains to be estimated. The absorption coefficient measures the probability of an interaction per frequency and unit distance. Given an Einstein coefficient X, it is given by 
\begin{align}
    \alpha(\nu) = \frac{h\nu}{c} X n \phi(\nu),
\end{align}
where $n$ is number density of atoms partaking in the interaction and $\phi$ is the line profile normalized to one. Then one finds for neutral hydrogen:
\begin{align}
    \alpha(\nu) 
    &= \frac{h\nu}{c}(B_{01}n_0 - B_{10}n_1)\phi(\nu)\\
    &\approx \frac{h\nu}{c}n_0 B_{01}\frac{h \nu}{k_B T_S}\phi(\nu),
\end{align}
the approximation coming from eq.~\ref{eq:Boltzmann}. The line profile is approximated by $\phi(\nu) \approx c/(s H(z) \nu)$ with $s$ a distance and $H$ the Hubble parameter at redshift z, an estimate made when considering the effect of cosmological recession. Then:
\begin{align}
    \tau(\nu) &= A(\nu) \frac{x_\text{HI}n_H }{T_S H(z)}\\
    A(\nu) &= \frac{3c^2 h A_{10}}{32 \pi \nu^2 k_B},
\end{align}
with $n_H$ the density of hydrogen atoms, for which one can use:
\begin{align}
    n_H = (1+\delta_b) \frac{3H_0^2}{8 \pi G m_p} \Omega_{b,0} (1-Y_\text{He})(1+z)^3,
\end{align}
with $\Omega_{b,0}$ the present-day baryon density, $m_p$ the proton mass, $G$ the gravitational constant, $\delta_b$ the fluctuation of baryon content and $Y_\text{He}$ the fraction in mass in helium.

Assuming matter domination yields $H(z) = H_0 \Omega^{1/2}_{m,0} (1 + z)^{3/2}$ with $H_0$ and $\Omega_{m,0}$ the present-day cosmological constant and matter density parameter, one obtains a more detailed approximation of the brightness offset temperature from~\ref{eq:delta_T_simple} in the limit $\tau \ll 1$:
\begin{align}
    \delta T_B \approx 27~x_{\text{HI}}(1+\delta_b)\biggl(1-\frac{T_\gamma}{T_S}\biggr)
    \biggl(\frac{1+z}{10}\frac{0.15}{\Omega_{m,0}h^2}\biggr)^{1/2}\biggl(\frac{\Omega_{b,0}h^2}{0.023}\biggr)~\si{mK},~\label{eq:brightness_offset_final}
\end{align} 
with $x_{\text{HI}}$ the neutral hydrogen fraction. From the expression in eq.~\ref{eq:brightness_offset_final}, it becomes clear that if $T_S > T_\gamma$, the measured brightness offset is positive and vice versa. 

\subsection{Qualitative behaviour of the 21cm signal}
During the Dark Ages, the expansion of the universe causes the neutral hydrogen gas to cool adiabatically, which causes the absorption signal for $z>30$. As collisions cease in the decreasingly dense gas, absorption stops and $\delta T_B \approx 0$. The emission of first Lyman-$\alpha$ photons causes $T_S < T_\gamma$, resulting in the strong absorption feature but soon thereafter saturating. These first sources also emit X-rays which heat the IGM, resulting in an emission feature. This signal decreases when X-rays have nearly completely heated the IGM, such that the signal is $\delta T_B \approx 0$ by $z\approx 6$ and remains so until today.
\begin{figure}[H]
    \centering
    \includegraphics[width=0.7\linewidth]{graphics/lightcone.png}
    \caption{A map of the $21~\si{cm}$ signal of neutral hydrogen~\cite{fundamentals}}
    \label{fig:lightcone}
\end{figure}

